\section*{Prefácio}

Se você está lendo este guia, parabéns: você decidiu assumir o controle da sua vida digital. 


Durante o nosso encontro no CMEI Araras, no dia 16 de março de 2026, compartilhamos histórias, medos e descobertas sobre como a tecnologia impacta o nosso dia a dia. Percebemos que a internet não precisa ser um lugar assustador e que não é necessário ser um gênio da informática para se proteger.


Muitas vezes, a tecnologia é desenhada para ser complexa, o que acaba afastando as pessoas e criando vulnerabilidades, principalmente em questões financeiras e na proteção das nossas crianças. A missão deste guia é quebrar essa barreira.


Assim como trancamos a porta de casa e olhamos para os dois lados antes de atravessar a rua, no mundo digital nós só precisamos criar novos (e simples) hábitos para navegar com paz interior.


Use este guia como o seu manual de sobrevivência. Configure seu celular no seu tempo, sem pressa. E, acima de tudo, compartilhe este conhecimento com suas amigas, alunos e familiares. A verdadeira segurança digital acontece quando cuidamos uns dos outros.


Uma excelente leitura e uma navegação segura!


\vspace{0.5cm}
\textbf{Naygno Barbosa Noia} \\
\textit{Ciência da Computação - UFBRA}

\vspace{1cm}
\begin{tcolorbox}[colback=NordBackground, colframe=BraveOrange, title=\faGithub\ Código Aberto (Open Source)]
	\textcolor{white}{Este guia foi totalmente tipografado em \LaTeX\ e seu código-fonte é aberto. Sinta-se à vontade para acessar o repositório no GitHub para sugerir melhorias ou compartilhar:}
	
	\vspace{0.3cm}
	\centering
	\href{https://github.com/naygno/guia-seguranca-cmei-araras}{\textbf{\textcolor{BitwardenBlue}{\faLink\ github.com/naygno/guia-seguranca-cmei-araras}}}
\end{tcolorbox}